% !TeX spellcheck = es_ES
\documentclass[letterpaper,11pt]{article}
\usepackage[utf8]{inputenc}
\usepackage[T1]{fontenc}
\usepackage[spanish]{babel}
\decimalpoint % separador decimal con punto, como en las tablas de la fuente original
\usepackage{lmodern,microtype}
\usepackage{amsmath,amssymb,amsfonts}
\usepackage{geometry}
\geometry{margin=2.54cm}
\usepackage{setspace}
\setstretch{1.15}
\usepackage{enumitem}
\usepackage{booktabs}
\usepackage{color}
\usepackage{hyperref}
\definecolor{ntr}{rgb}{0.8,0,0}
\hypersetup{colorlinks=true,linkcolor=ntr,citecolor=ntr,urlcolor=black}

% Portada (Luis Cabral): titulo \LARGE sans-serif alineado a la izquierda, sin numero de pagina
\makeatletter
    \renewcommand{\@maketitle}{%
    \newpage\parindent0in\null\vskip2em%
    {\LARGE\textbf{\textsf{\@title}}\par\vskip2em}%
    \@author\par\vskip2em%
    \@date%
    \par
    \thispagestyle{empty}\setcounter{page}{1}%
    }%
\makeatother

\DeclareMathOperator{\Var}{Var}
\DeclareMathOperator{\Cov}{Cov}
\DeclareMathOperator{\Corr}{Corr}
\DeclareMathOperator{\E}{\mathbb{E}}
\DeclareMathOperator{\N}{Normal}
\DeclareMathOperator{\LP}{\mathbb{L}}
\newcommand\independent{\protect\mathpalette{\protect\independenT}{\perp}}
\def\independenT#1#2{\mathrel{\rlap{$#1#2$}\mkern2mu{#1#2}}}
% Teoria asintotica
\newcommand{\SR}{\mathbb{R}}
\DeclareMathOperator{\plim}{plim}
\DeclareMathOperator{\Avar}{AVar}
\newcommand{\pto}{\stackrel{p}{\longrightarrow}}
\newcommand{\dto}{\stackrel{d}{\longrightarrow}}
\newcommand{\adistr}{\stackrel{a}{\sim}}
\newcommand{\epsi}{\varepsilon}


\title{Análisis Econométrico \\ Ayudantía 7}
\author{Magíster en Economía, Universidad Alberto Hurtado}
\date{}

\begin{document}
\maketitle

\noindent Esta ayudantía comienza con un control breve (\textbf{Control 3}, 15 minutos, sin
apuntes) sobre el contenido de la Tarea 3. Luego revisamos en conjunto los ejercicios 1, 3, 4
y 5 de la Tarea 3 --- se reproducen aquí para tenerlos a mano durante la revisión. Los datos
son \texttt{wage2.rds} y \texttt{hprice2.rds}, en la carpeta \texttt{problem-sets/datos/} del
repositorio del curso.

\vspace{1em}

\noindent\textbf{Ejercicio 1 de la Tarea 3 (20 puntos). Varianza asintótica bajo
homoscedasticidad.}

Suponga homoscedasticidad. En clases vimos que, en ese caso, la varianza asintótica del
estimador MCO es $\Avar(\hat \beta)=\frac{\sigma^2}{N}\E(x_ix_i')^{-1}$.
\begin{enumerate}
    \item Considere el modelo de regresión simple
    \begin{align*}
    y_i=\beta_0+\beta_1 x_{1i}+u_i,
    \end{align*}
    donde $\E(u_i)=0$, $\E(x_{1i}\, u_i)=0$ y $\E(u_i^2|x_{1i})=\sigma^2$. Utilice la fórmula
    anterior con $x_i=(1, x_{1i})'$ para demostrar que
    \begin{align*}
    \Avar(\hat \beta_1)=\frac{\sigma^2}{N}\frac{1}{\Var(x_{1i})}.
    \end{align*}
    \item Considere ahora el modelo con dos variables explicativas
    \begin{align*}
    y_i=\beta_0+\beta_1 x_{1i}+\beta_2 x_{2i}+v_i,
    \end{align*}
    donde $\E(v_i)=\E(x_{1i}\,v_i)=\E(x_{2i}\,v_i)=0$ y $\E(v_i^2|x_{1i},x_{2i})=\sigma_v^2$. En
    clases vimos que
    \begin{align*}
    \Avar(\hat \beta_1)=\frac{\sigma_v^2}{N}\frac{1}{\Var(r_{1i})},
    \end{align*}
    donde $\gamma_0+\gamma_1x_{2i}$ es el proyector lineal de $x_{1i}$ en una constante y
    $x_{2i}$, y $r_{1i}=x_{1i}-\gamma_0-\gamma_1x_{2i}$ es su residuo. Demuestre que
    \begin{align*}
    \Var(r_{1i})=\Var(x_{1i})\,(1-\rho^2),
    \end{align*}
    donde $\rho=\Corr(x_{1i},x_{2i})$ es el coeficiente de correlación entre $x_{1i}$ y
    $x_{2i}$.
    \item Explique cómo depende $\Avar(\hat \beta_1)$ del tamaño de la muestra, de la varianza
    del error, de la varianza de $x_1$ y de la correlación entre $x_1$ y $x_2$. ¿Qué ocurre
    cuando $|\rho|\to 1$?
\end{enumerate}

\vspace{1em}

\noindent\textbf{Ejercicio 3 de la Tarea 3 (20 puntos). Dummies e interacciones.}

Utilice los datos de \texttt{wage2} del trabajo de Blackburn y Neumark (QJE, 1992): 935
hombres de entre 28 y 38 años en 1980. Las variables \texttt{wage} (salario mensual, en
dólares), \texttt{married} (igual a 1 si está casado) y \texttt{urban} (igual a 1 si vive en
una zona urbana) definen el modelo
\begin{align*}
\log(wage_i)=\beta_0+\beta_1 \, married_i+\beta_2 \, urban_i+\beta_3 \, married_i\times urban_i+u_i.
\end{align*}
\begin{enumerate}
    \item Estime el modelo. Escriba $\E(\log(wage)|married,urban)$ para los cuatro grupos
    (casado urbano, casado rural, soltero urbano y soltero rural) en función de los $\beta$'s,
    e interprete cada uno de los coeficientes estimados.
    \item Calcule la media muestral de $\log(wage)$ en cada grupo y verifique que coincide con
    el valor predicho por el modelo. Explique por qué.
    \item Suponga que construye una dummy para cada uno de los cuatro grupos y estima la
    regresión de $\log(wage)$ en las cuatro dummies y una constante. ¿Qué ocurre? Explique.
\end{enumerate}

\vspace{1em}

\noindent\textbf{Ejercicio 4 de la Tarea 3 (30 puntos). Retornos de la educación.}

Utilice los datos de \texttt{wage2} para estudiar los retornos de la educación. Además de las
variables del ejercicio anterior, la base incluye \texttt{educ} (años de educación),
\texttt{IQ} (resultado en un test de inteligencia), \texttt{exper} (años de experiencia
laboral), \texttt{tenure} (años en el empleo actual) y \texttt{south} (igual a 1 si vive en el
sur de EE.UU.).
\begin{enumerate}
    \item Utilice un modelo de regresión simple log-nivel para estimar el efecto de un año más
    de educación ($educ$) en salarios ($wage$). Interprete los resultados. ¿Usted piensa que el
    coeficiente estimado mide una relación causal? Explique.
    \item El modelo de regresión simple omite, entre otras cosas, la habilidad innata. En los
    datos existe una proxy de la habilidad innata: $IQ$. Antes de estimar, prediga si el
    coeficiente de $educ$ aumentará o disminuirá al incluir $IQ$ como variable explicativa
    adicional, y explique por qué. Luego estime el modelo con $IQ$, interprete el coeficiente
    de $educ$ y compare con (a).
    \item La fórmula del sesgo de variable omitida también se cumple en la muestra,
    reemplazando momentos poblacionales por momentos muestrales. Primero, compruebe este
    resultado en los datos: regrese $IQ$ en $educ$ y, en base a esa estimación y a lo obtenido
    en (b), recupere la pendiente estimada en (a). Segundo, demuestre analíticamente que este
    resultado se cumple siempre en la muestra.
    \item Estime el efecto de la educación sobre los salarios controlando por $IQ$, $exper$,
    $tenure$, $married$, $south$ y $urban$. Compare con los resultados de (a) y (b).
\end{enumerate}

\vspace{1em}

\noindent\textbf{Ejercicio 5 de la Tarea 3 (20 puntos). Contaminación y precio de las
viviendas.}

Utilice los datos de \texttt{hprice2} del trabajo de Harrison y Rubinfeld (JEEM, 1978) sobre
el efecto de la contaminación en el precio de las viviendas: 506 comunidades del área
metropolitana de Boston en 1970. Las variables son \texttt{price} (precio mediano de las
viviendas, en dólares), \texttt{nox} (concentración de óxido de nitrógeno en el aire, en
partes por 100 millones), \texttt{rooms} (número promedio de habitaciones), \texttt{dist}
(distancia ponderada a cinco centros de empleo) y \texttt{stratio} (ratio estudiantes por
profesor).
\begin{enumerate}
    \item Utilice un modelo de regresión simple log-log para estimar el efecto de un aumento en
    $nox$ en $price$. Interprete los resultados. ¿Usted piensa que el coeficiente estimado mide
    una relación causal? Explique.
    \item Agregue las siguientes variables explicativas al modelo: $rooms$, $\log(dist)$ y
    $stratio$. Interprete la estimación del efecto de un aumento de la contaminación en el
    precio de las viviendas, compare con (a) y explique la diferencia.
    \item Obtenga la estimación del coeficiente de $\log(nox)$ del modelo en (b) utilizando la
    regresión particionada. Primero, guarde los residuos de una regresión de $\log(nox)$ en
    $rooms$, $\log(dist)$ y $stratio$. Segundo, regrese $\log(price)$ en esos residuos.
    Verifique que la estimación coincide con (b). No se preocupe por los errores estándar,
    sólo por las estimaciones puntuales.
    \item Compare el error estándar del coeficiente de $\log(nox)$ en (a) y en (b). Calcule la
    correlación muestral entre $\log(nox)$ y $\log(dist)$ y utilice la intuición del ejercicio 1
    para explicar la diferencia.
\end{enumerate}

\vspace{2em}
\noindent\rule{\textwidth}{0.4pt}

\section*{Material adicional (opcional)}

\noindent Para discutir si queda tiempo.

\vspace{1em}

\noindent\textbf{Ejercicio 2 de la Tarea 3 (10 puntos). Horas de sueño y mortalidad.}

Un estudio reciente ha encontrado que la tasa de mortalidad de las personas que duermen 7 horas
por noche es menor que la tasa de mortalidad de aquellas personas que duermen 8 o más horas. El
estudio utilizó 1,1 millones de observaciones de una muestra aleatoria de chilenos entre 30 y
102 años de edad. Cada persona en la muestra fue seguida por 4 años. La tasa de mortalidad para
personas que duermen 7 horas fue calculada como el número de muertes durante los 4 años entre
las personas que duermen 7 horas sobre el número de personas que duermen 7 horas en la muestra.
Cálculos similares se realizaron para el resto de las categorías. En base a los resultados del
estudio, ¿usted recomendaría a los chilenos que duermen 9 horas por noche reducir las horas de
sueño a 7 horas por noche si quieren prolongar sus vidas? ¿Por qué o por qué no? Explique.

\end{document}
